\documentclass[11pt]{article}
\usepackage[margin=1in]{geometry}
\usepackage[T1]{fontenc}
\usepackage[utf8]{inputenc}
\usepackage{lmodern}
\usepackage{microtype}
\usepackage{booktabs}
\usepackage{array}
\usepackage{enumitem}
\usepackage{hyperref}
\setlength{\parindent}{0pt}
\setlength{\parskip}{0.65em}

\title{Prompt-Profile Objective: Plain-Language Note}
\author{}
\date{2026-03-30 09:14 UTC}

\begin{document}
\maketitle

\section*{Bottom line}
The project wants a predictor that looks only at prompt-prefill activations and
flags prompts that are likely to make the model go bad under one fixed model and
one fixed decode policy. On the current saved five-dataset bundle, the best main
target for that job is \texttt{p\_loop}.

\section*{What the project is actually trying to do}
We are not trying to predict whether a prompt is ``inherently loopy'' in the
abstract. We are trying to predict what usually happens for one fixed
\texttt{(model, prompt, decode policy)} tuple before generation starts.

The practical use is prompt screening: rank prompts by risk, then inspect or
avoid the worst ones. ``Bad'' here mainly means:
\begin{itemize}[leftmargin=1.5em]
\item the rollout falls into the loop regime;
\item the rollout burns a lot of budget and may hit the cap;
\item accuracy drops sharply inside that risky slice.
\end{itemize}

\section*{Definitions}
\textbf{\texttt{p\_loop}}: for one prompt, how often repeated rollouts loop. If
\texttt{3} out of \texttt{10} rollouts loop, then \texttt{p\_loop = 0.3}.

\textbf{\texttt{p\_cap}}: for one prompt, how often repeated rollouts hit the
hard max-length cap.

\textbf{\texttt{mean\_relative\_length}}: on average, what fraction of the
available generation budget the prompt consumes.

\textbf{\texttt{majority\_s\_0.5}}: a binary label that is \texttt{1} if more
than half of repeated rollouts use at least half of the available budget. This
is better read as ``often long'' than as ``often loops.''

\textbf{\texttt{effective\_budget}}: the actual generation budget available
after prompt length and context-limit constraints are taken into account.

\textbf{Loop-rate enrichment}: loop rate inside the top predicted-risk bucket,
divided by the overall loop rate on the test split. A value of \texttt{2.0x}
means the flagged bucket has twice the average loop rate.

\section*{What was missing from the last thread}
The last thread was still missing two decisive pieces:
\begin{enumerate}[leftmargin=1.5em]
\item real metadata-only baselines for the continuous heads, not just raw
\texttt{Spearman(prompt\_length, target)};
\item an operational bucket test: rank held-out prompts by predicted risk, keep
the top \texttt{20\%}, and measure the actual loop rate, cap-hit rate, and
accuracy inside that bucket.
\end{enumerate}

Without those two pieces, the old thread could not cleanly answer which
objective should actually be trained.

\section*{What was trained and compared}
This pass used the finished saved prompt-profile bundles for:
\texttt{GPQA}, \texttt{AIME}, \texttt{MATH-500}, \texttt{MMLU-Pro}, and
\texttt{LiveCodeBench}.

The comparison surface contained:
\begin{itemize}[leftmargin=1.5em]
\item existing activation probes already trained from prompt-prefill activations
for \texttt{majority\_s\_0.5}, \texttt{p\_loop}, and
\texttt{mean\_relative\_length};
\item two existing activation readouts: last-layer MLP and per-layer ensemble
MLP;
\item new metadata-only controls trained in this pass:
  \begin{itemize}[leftmargin=1.5em]
  \item logistic regression for \texttt{p\_loop};
  \item linear regression for \texttt{mean\_relative\_length};
  \item features = \texttt{prompt\_length}, \texttt{effective\_budget}, and
  both together.
  \end{itemize}
\end{itemize}

The split stayed prompt-disjoint: train on one set of prompts, test on a
different set, and never let the same prompt leak across the split through
repeated rollouts.

\section*{How the comparison worked}
For each prompt:
\begin{enumerate}[leftmargin=1.5em]
\item run repeated rollouts under one fixed decode policy;
\item turn those rollouts into prompt-level target numbers such as
\texttt{p\_loop};
\item train on train prompts only;
\item score held-out test prompts;
\item sort test prompts by predicted risk;
\item keep the top \texttt{20\%} highest-risk prompts;
\item measure what actually happened inside that bucket:
  \begin{itemize}[leftmargin=1.5em]
  \item true loop rate;
  \item true cap-hit rate;
  \item true accuracy when prompt-level correctness exists.
  \end{itemize}
\end{enumerate}

This bucket test is the key operational check. It answers: if we really use this
score to flag bad prompts, does it isolate the bad slice?

\section*{Why \texttt{p\_loop} wins}
\begin{center}
\begin{tabular}{>{\raggedright\arraybackslash}p{1.55in}ccc}
\toprule
Dataset & \texttt{majority\_s\_0.5} & \texttt{p\_loop} & strongest metadata \\
 & loop enrich. & loop enrich. & baseline \\
\midrule
\texttt{GPQA} & \texttt{1.82x} & \texttt{1.82x} & \texttt{0.91x} \\
\texttt{AIME} & \texttt{0.80x} & \texttt{3.20x} & \texttt{0.80x} \\
\texttt{MATH-500} & \texttt{2.50x} & \texttt{3.50x} & \texttt{1.50x} \\
\texttt{MMLU-Pro} & \texttt{2.61x} & \texttt{3.48x} & \texttt{1.74x} \\
\texttt{LiveCodeBench} & \texttt{1.63x} & \texttt{2.32x} & \texttt{1.26x} \\
\bottomrule
\end{tabular}
\end{center}

Average loop-rate enrichment across datasets:
\begin{itemize}[leftmargin=1.5em]
\item \texttt{p\_loop}: \texttt{2.86x}
\item \texttt{majority\_s\_0.5}: \texttt{1.87x}
\item \texttt{mean\_relative\_length}: \texttt{1.77x}
\item strongest metadata baseline: \texttt{1.24x}
\end{itemize}

On the four datasets where prompt-level accuracy is available
(\texttt{GPQA}, \texttt{AIME}, \texttt{MATH-500}, \texttt{MMLU-Pro}), the
\texttt{p\_loop} top-risk bucket is also the lowest-accuracy bucket every time.

If the operational goal is ``find prompts that actually go bad,'' then
\texttt{p\_loop} is the target most aligned with that job: it directly asks how
often the prompt loops, it consistently concentrates loop-heavy prompts, and it
also concentrates low-accuracy prompts where accuracy is available.

\section*{Why the other objectives lost}
\textbf{Why not \texttt{majority\_s\_0.5}}: it is a coarse thresholded label. It
mainly asks whether the prompt is often long, not whether it often loops.
Prompt length alone already does very well on this label in some slices,
especially \texttt{AIME}. So it still has signal, but it is too geometry-heavy
to be the main headline objective.

\textbf{Why not \texttt{mean\_relative\_length}}: it is a useful dense target and
ranks budget usage well, but it mixes correct long reasoning, wrong long
reasoning, and looping. It is good as a secondary utility or budget head, not
the cleanest main bad-prompt screen.

\textbf{Why not \texttt{p\_cap}}: it is too narrow. Many bad looped rollouts do
not hit the cap, so training only on cap hits would miss a lot of the failure
mass we care about.

\section*{Why this training objective leads to the best result}
The main reason is alignment between the target and the downstream decision:
\begin{itemize}[leftmargin=1.5em]
\item training \texttt{majority\_s\_0.5} teaches a coarse ``often long'' label;
\item training \texttt{mean\_relative\_length} teaches budget usage;
\item training \texttt{p\_cap} teaches only the narrow cap-hit subset;
\item training \texttt{p\_loop} teaches the event we actually want to screen
for: entering the loop regime.
\end{itemize}

So the winning result is not mysterious. The objective itself matches the
operational question better, and that advantage survives the held-out bucket
test.

\section*{What this resolves, and what it does not}
This pass resolves the main open issue from the last thread:
\begin{itemize}[leftmargin=1.5em]
\item the decision is no longer resting on one dataset anecdote;
\item it is no longer resting on raw prompt-length correlation for the
continuous heads;
\item it is no longer missing the operational bucket test.
\end{itemize}

So the target choice is now materially better supported:
\begin{itemize}[leftmargin=1.5em]
\item main objective: \texttt{p\_loop}
\item secondary utility head: \texttt{mean\_relative\_length}
\item control / cheap auxiliary screen: \texttt{majority\_s\_0.5}
\item do not reopen direct \texttt{p\_cap} now
\end{itemize}

What it still does \textbf{not} resolve:
\begin{enumerate}[leftmargin=1.5em]
\item this is retrospective on saved bundles, not a fresh prospective run;
\item \texttt{LiveCodeBench} still lacks prompt-level accuracy in the recovered
projection artifact, so the low-accuracy claim is \texttt{4 / 5} datasets;
\item small datasets, especially \texttt{AIME}, can still move noticeably when
only a few prompts change buckets.
\end{enumerate}

\section*{Exact recommended next run}
\begin{itemize}[leftmargin=1.5em]
\item Freeze the architecture, feature view, and checkpoint-selection rule up
front.
\item Train one fresh prompt-disjoint \texttt{p\_loop} predictor on a
representative hard slice.
\item Keep the metadata baselines mandatory:
  \texttt{prompt\_length}, \texttt{effective\_budget}, and both together.
\item Train \texttt{mean\_relative\_length} only if a secondary utility or
budget score is still wanted.
\item Keep \texttt{majority\_s\_0.5} in the evaluation bundle as a control, not
as the main target.
\end{itemize}

\section*{Bottom line again}
The plainest summary is:
\begin{itemize}[leftmargin=1.5em]
\item the project goal is to flag prompts that make the model go bad before
generation starts;
\item the missing baselines and bucket test are now done on the saved
five-dataset bundle;
\item \texttt{p\_loop} is the objective most aligned with that goal and the one
that wins the held-out risk-bucket test;
\item \texttt{mean\_relative\_length} is still useful, but as a secondary budget
head;
\item \texttt{majority\_s\_0.5} still has signal, but it is too close to a
geometry-heavy ``often long'' control to be the main objective.
\end{itemize}

\end{document}
